\section{The actual pro-\texorpdfstring{$\ell$}{ell} quotient and matched descent}
\label{sec:proell-descent-v039}

\status{PQL1--PQL21 reconstruct the original quotient system and the
implication \(\mathrm{CoreBC}\Rightarrow\mathrm{Raw}\). All descended
points retain their coefficient field and their original map out of
\(B\); every automorphic module and the ideal \(J\) are unchanged.
The normalized Core producer and its patched support are separate
obligations.}

The formal source locators accompanying this chapter refer to the
fixed commit
\par\smallskip\noindent
\texttt{7cbf7a0c32fca1badf05931879a37c3f8a2b03b4}.\par\smallskip
Correctness of that formal proof is the standing assumption.
The three determinant results used from the literature are isolated
in PQL7, with their actual hypotheses checked.

% BEGIN PQL1
\paragraph{PQL1. The exact descent problem.}\label{claim:PQL1}
Let \(\ell\) be prime, \(L\) the original finite list of local
profinite groups, and \(G_0\) the original global group with its
sign character \(\eta\), involution \(c\), cyclotomic character
\(\varepsilon\), and all local maps. A morphism of such groups
is a continuous surjection preserving every one of these data.
Write \(\Gamma^+=\ker\eta\). Let \(T\) be the given
\(\mathrm{TWGrp}\), with open normal subgroup \(H\subset\Gamma_0\),
surjections \(h_N:\Gamma_N\to\Gamma_0\), and \(q=T.q\) pairs
\((\phi_{N,j},\tau_{N,j})\). Put \(H_N=h_N^{-1}H\).
The definition includes
\[
 H\subset\ker\eta,\qquad \varepsilon(H)\subset1+\ell\mathbb Z_\ell,
 \qquad |\mathrm{Hom}_{\mathrm{cont}}(H_N,\mathbb F_\ell)|
       \leq T.hb,                                      \tag{PQL.1}
\]
as well as the original kernel, tame-relation and convergence clauses.
Here additive \(\mathbb F_\ell\) is the source's multiplicative
\(\mathrm{ZMod}\,\ell\) notation, not a new coefficient field.

For any open normal \(H\) in a profinite group \(\Gamma\), define
\[
 P_\ell^\Gamma(H)=
 \bigcap_{\substack{V\triangleleft_{\!o}\Gamma\\
              H/(H\cap V)\ {\rm an}\ \ell{\rm\!-\!group}}}V.
                                                            \tag{PQL.2}
\]
This is exactly the source's \(\mathrm{PL}\,\ell\,H\):
in a finite group, every element having \(\ell\)-power order is
equivalent to the group being an \(\ell\)-group, by Cauchy's theorem.
The task is to quotient simultaneously by \(P_\ell^{\Gamma_0}(H)\)
and \(P_\ell^{\Gamma_N}(H_N)\), and to transport all arithmetic
data through these very quotient maps.

For the arithmetic part retain the original coefficient DVR \(O\),
complete local Noetherian \(O\)-algebra \(B\), residual datum \(r\),
determinant pair \(u=(D,\psi)\), ideal \(J\), and
\(\mathrm{TWAut}(u,J,\sigma,D_0,T)\).
Here \(D_0\) denotes the source's weight-variable number; below \(d\)
will denote the bound used in \(\mathrm{Pt}\).
The source's \(\mathrm{Adapted}\) condition says that all residual
characters are trivial on \(H\), and that \(r.Ip_w=G_0.Iu_w\).
No residual absolute irreducibility is assumed. In the final
application \(\sigma=\mathrm{Empty}\), \(n=4\), and \(\ell=3\).
The proof also preserves an arbitrary finite chain of original
points, so its conclusion is the full source implication
\(\mathrm{CoreBC}\Rightarrow\mathrm{Raw}\), not merely its
one-point specialization.
% END PQL1

% BEGIN PQL2
\paragraph{PQL2. The relative pro-\(\ell\) residual.}\label{claim:PQL2}
For a profinite group \(H\) let
\[
 R_\ell(H)=\bigcap_{\substack{U\triangleleft_{\!o}H\\H/U\
                                      {\rm an}\ \ell{\rm\!-\!group}}}U .
 \qquad
 P_\ell^\Gamma(H)=R_\ell(H)\subset H
       \quad(H\triangleleft_o\Gamma).                \tag{PQL.3}
\]
The intrinsic residual is closed and characteristic in \(H\).
To prove the equality, intersecting any admissible \(V\) in
(PQL.2) with \(H\) shows \(R_\ell(H)\subset P_\ell^\Gamma(H)\);
the admissible choice \(V=H\) puts the latter subgroup inside \(H\).
Conversely take \(U\triangleleft_oH\) with \(H/U\) an
\(\ell\)-group. Its core \(V=\bigcap_{\gamma\in\Gamma}
\gamma U\gamma^{-1}\) is open normal in \(\Gamma\).
Only finitely many conjugates occur because \(U\) is open in
\(\Gamma\). Each \(H/\gamma U\gamma^{-1}\) is an
\(\ell\)-group, and \(H/V\) embeds in their finite product.
Thus \(V\) is admissible in (PQL.2), and \(P_\ell^\Gamma(H)
\subset V\subset U\), proving the equality and normality.

The quotient \(H/R_\ell(H)\) is pro-\(\ell\): its natural map
into the product of the finite quotients \(H/U\) is a continuous
injection with compact, hence closed, image. Conversely a
pro-\(\ell\) group has trivial residual, since its finite
\(\ell\)-quotients separate points. We obtain
\[
 P_\ell^\Gamma(H)=1\quad\Longleftrightarrow\quad
 H\ \hbox{is pro-}\ell.                              \tag{PQL.4}
\]
For \(V\triangleleft_o\Gamma\) we also have
\[
 P_\ell^\Gamma(H)\subset V
 \quad\Longleftrightarrow\quad H/(H\cap V)
                         \ \hbox{is an }\ell\hbox{-group}.   \tag{PQL.5}
\]
In the forward direction the quotient is a finite continuous
image of \(H/R_\ell(H)\); the reverse direction follows from
the defining intersection. This proves the finite-quotient
content of the original \texttt{pl\_props}.
% END PQL2

% BEGIN PQL3
\paragraph{PQL3. Functoriality and passage to an open subgroup.}
\label{claim:PQL3}
If \(f:\Gamma\to\Gamma'\) is a continuous surjection and
\(H'\triangleleft_o\Gamma'\), put \(H=f^{-1}H'\). Then
\[
 f(P_\ell^\Gamma(H))=P_\ell^{\Gamma'}(H').             \tag{PQL.6}
\]
The inclusion from left to right follows by pulling back finite
\(\ell\)-quotients of \(H'\). The image on the left is closed,
and \(H'/f(R_\ell(H))\) is a continuous image of
\(H/R_\ell(H)\), hence pro-\(\ell\). The kernels of its
finite quotients, pulled back to \(H'\), show that
\(R_\ell(H')\subset f(R_\ell(H))\).

We also need the source's \texttt{pl\_sub} in the following form.
Suppose \(H\triangleleft\Gamma\), \(K\triangleleft_o\Gamma\),
\(H\subset K\), and \(U\triangleleft_oK\). If every element
of the image of \(H\) in \(K/U\) has \(\ell\)-power order, then
\[
 P_\ell^\Gamma(H)\cap K\subset U.                    \tag{PQL.7}
\]
In our applications \(H\) is open. Take the core \(V\) of
\(U\) in \(\Gamma\). The finite image of the normal subgroup
\(H\) in each \(K/\gamma U\gamma^{-1}\) is an \(\ell\)-group;
normality of \(H\) identifies these images with the image
in \(K/U\). Its image modulo their finite intersection is
therefore an \(\ell\)-group. Equation (PQL.2) gives
\(P_\ell^\Gamma(H)\subset V\subset U\).
This keeps the ambient group in the definition of
\(\mathrm{PL}\) distinct from its open subgroup \(\Gamma^+\).
% END PQL3

% BEGIN PQL4
\paragraph{PQL4. The quotient preserves the global frame.}
\label{claim:PQL4}
Set \(P=P_\ell^\Gamma(H)\), where \(H\subset\ker\eta\) and
\(\varepsilon(H)\subset1+\ell\mathbb Z_\ell\).
The quotient \(\Gamma/P\) is profinite: closedness of \(P\)
gives a compact Hausdorff quotient, and open normal subgroups
of \(\Gamma\) containing \(P\) separate its points.

Both characters descend. For \(\eta\) use \(P\subset H\).
For \(\varepsilon\), reduce modulo \(\ell^a\).
The image of \(1+\ell\mathbb Z_\ell\) is an \(\ell\)-group:
each successive kernel from modulus \(\ell^b\) to
\(\ell^{b-1}\) embeds additively in \(\mathbb F_\ell\).
By (PQL.5), \(P\) lies in the kernel of every reduction.
Separation by these reductions gives \(\varepsilon(P)=1\).
This argument also works when \(\ell=2\).

Take \(c'\) to be the image of \(c\), and compose every
original local map with the quotient. The relations
\((c')^2=1\), \(\eta'(c')=-1\), and all character compatibilities
follow by composition. Continuity of the descended
cyclotomic character follows from the quotient topology;
the descended sign kernel is open because its preimage is
the original open kernel. Thus the quotient map is exactly
a \(\mathrm{Glob.Hom}\). This proves \texttt{quot\_ex},
with the original local list \(L\).
% END PQL4

% BEGIN PQL5
\paragraph{PQL5. All fields of the original \(\mathrm{TWQuot}\).}
\label{claim:PQL5}
Apply PQL4 to \(H\subset\Gamma_0\) and \(H_N\subset\Gamma_N\),
and write the quotient maps \(g_0,g_N\). By (PQL.6),
\(h_N(P_\ell^{\Gamma_N}(H_N))=P_\ell^{\Gamma_0}(H)\);
there is a unique continuous surjection \(h'_N\) with
\[
 h'_Ng_N=g_0h_N.                                    \tag{PQL.8}
\]
It preserves the frames by surjectivity. Put \(H'=g_0(H)\).
Quotient maps of groups are open; since \(H\) contains
\(\ker g_0\), this is open normal and \(g_0^{-1}H'=H\).
Similarly
\[
 g_N^{-1}((h'_N)^{-1}H')=H_N.                        \tag{PQL.9}
\]
To verify convergence of the \(\phi_{N,j}\), pull an open
normal subgroup of \(\Gamma'_0\) back along \(g_0\).
It contains \(P_\ell^{\Gamma_0}(H)\), so the original
convergence clause applies. The elements \(\sigma_j\),
\(\phi_{N,j}\), and \(\tau_{N,j}\) become their literal
images. The integers \(q,Nq(N,j)\) and their congruences
are unchanged. The tame relation becomes an equality,
because its original error lies in \(\ker g_N\).

For the kernel condition, lift \(x'\in\ker h'_N\) to
\(x\in\Gamma_N\). Then \(h_N(x)\in\ker g_0=h_N(\ker g_N)\).
Choose \(y\in\ker g_N\) with \(h_N(y)=h_N(x)\). We have
\(xy^{-1}\in\ker h_N\) and \(g_N(xy^{-1})=x'\).
The original kernel clause puts \(xy^{-1}\) in the closure
of the subgroup normally generated by the \(\tau_{N,j}\)
and \(\ker g_N\). Its image lies in the closed normal
subgroup generated by the \(g_N(\tau_{N,j})\), which is
exactly the required descended clause.

Pullback injects continuous homomorphisms from \(H'\) and
\((h'_N)^{-1}H'\) to \(\mathbb F_\ell\) into the original
\(\mathrm{HomL}\) sets. It preserves open kernels, and
surjectivity gives injectivity. Thus finiteness and the
same numerical bound \(T.hb\) survive. Finally (PQL.6)
for \(g_0,g_N\), together with (PQL.9), makes both
descended \(\mathrm{PL}\) subgroups trivial.
The \(\mathrm{PhiCond}\) predicates refer only to the
unchanged local groups and are retained verbatim.
Every field of \(\mathrm{TWQuot}\) and its associated
\(Q.\mathrm{grp}\) is now checked:
\[
 \mathrm{Nonempty}(\mathrm{TWQuot}\,T),
 \qquad (Q.\mathrm{grp}).\mathrm{ProL}.              \tag{PQL.10}
\]
No limit system or arithmetic point was used.
% END PQL5

% BEGIN PQL6
\paragraph{PQL6. Residual characters and local data descend.}
\label{claim:PQL6}
Let \(g:G\to G'\) be a quotient morphism with
\(\ker g=P_\ell^\Gamma(H)\). Suppose the residual characters
\(\bar\chi_i\) are trivial on \(H\). They factor uniquely
to characters \(\bar\chi'_i\) on \((\Gamma')^+\), because
\(\ker g\subset H\). Their kernels are open by the quotient
topology, and pairwise distinctness follows by surjective
pullback. Keep \(\bar\psi\) and the integral local characters
\(\theta\) literally unchanged.

For every commutative coefficient algebra \(C\), the
group-algebra map \(p_*:C[\Gamma^+]\to C[(\Gamma')^+]\)
is surjective. It commutes with polarization, which is
formed from \(c\), inversion and the cyclotomic factor,
all preserved by \(g\). The diagonal residual
representations agree after pullback, so their full
determinant laws agree there too. Lifting any group-algebra
element proves residual polarization downstairs.
This verifies \(\mathrm{Resid}\) and
\(\mathrm{Glob.Compat}(g,r,r')\).

If \(r.Ip_w=G.Iu_w\), all residual characters are trivial
on the original inertia image. The equality
\(g\phi_u=\phi'_u\) gives triviality on the new image,
so \(r'.Ip_w=G'.Iu_w\). Thus the local part of
\(\mathrm{Adapted}\) also descends. No new local
character or weight is selected.
% END PQL6

% BEGIN PQL7
\paragraph{PQL7. Three checked determinant inputs.}\label{claim:PQL7}
We use three results for multiplicative homogeneous polynomial
laws, with kernel always referring to the full law over every
commutative scalar algebra. Chenevier's Lemma 1.19 identifies it as
\[
 \{z:\ D_C(1+hz)=1\ \hbox{for every }C,h\}
\]
and gives factorization through a surjective algebra quotient
whose kernel it contains. Lemma 1.21 places the coefficient
Cayley--Hamilton ideal in this kernel. Lemma 2.33 says that a
continuous determinant on a profinite group with discrete
coefficients kills the augmentation ideal of some open
normal subgroup.\footnote{G. Chenevier,
\emph{The \(p\)-adic analytic space of pseudocharacters of a
profinite group and pseudorepresentations over arbitrary rings},
Lemmas 1.19, 1.21 and 2.33, printed pages 17--18 and 38.
\url{https://maths.dur.ac.uk/events/Meetings/LMS/2011/GRAF11/papers/chenevier.pdf}.}

The source's \(\mathrm{GDet}\) is such a law on the group
algebra. Its \(\mathrm{CH}\) is the set of coefficients of
characteristic-polynomial identities for generic finite
linear combinations, so its generated ideal is the ideal
of Lemma 1.21. At coefficient level \(B/\mathfrak m_B^a\),
\(\mathrm{IsCont}\) makes each characteristic coefficient
locally constant, exactly the condition in Lemma 2.33.
The group-algebra quotient is surjective after every
commutative scalar extension, as needed in Lemma 1.19.
There is no reducedness, residual irreducibility, or
invertibility-of-\(n!\) hypothesis. If \(n=0\), the law
is identically one and these conclusions are immediate;
the points used in Raw have \(n>0\), since \(0<a\leq n\).
% END PQL7

% BEGIN PQL8
\paragraph{PQL8. Closedness of the full determinant kernel.}
\label{claim:PQL8}
Let \(I\) be an ideal of an \(I\)-adically separated ring
\(B\), and \(D\) an \(I\)-adically continuous determinant
on a profinite group \(G\). For \(S\subset B[G]\), set
\[
 \overline{(S)}^{\,I,G}
 =\bigcap_{a\geq0,\ V\triangleleft_oG}
       (S,\ I^a,\ [v]-1\ (v\in V)).                 \tag{PQL.11}
\]
Each term is a two-sided ideal. This is precisely
\(\mathrm{closureIdeal}\ I\ S\).
If \(D\) kills every element of \(S\), it kills this intersection.

For each \(a\), PQL7 supplies an open normal subgroup
\(V_a\) whose augmentation ideal is killed by the reduced
law \(D_a\) over \(B/I^a\). Its kernel is a two-sided
ideal containing the images of \(S\) and \(I^a\), so the
image of any \(z\) in (PQL.11) is in \(\ker D_a\).
For arbitrary finitely many \(h_j\in B[G]\), form
\[
 D_{B[X_1,\ldots,X_s]}
       \left(1+\left(\sum_j X_jh_j\right)z\right)-1 .
                                                            \tag{PQL.12}
\]
Every coefficient vanishes modulo every \(I^a\), and
therefore is zero. Substituting coefficients in any
commutative \(B\)-algebra proves \(D_C(1+hz)=1\)
for every \(C,h\), since each \(h\) is a finite sum.
This proves closedness for the full law, not just its trace.
The original \(B\) is separated by \(\mathrm{IsCO}\).
Taking \(S=\mathrm{CH}(D)\) gives
\[
 \overline{(\mathrm{CH}(D))}^{\,\mathfrak m_B,G}
                                      \subset\ker D.        \tag{PQL.13}
\]
% END PQL8

% BEGIN PQL9
\paragraph{PQL9. Killing the residual in every closed Cayley--Hamilton ideal.}
\label{claim:PQL9}
Suppose \(u\) reduces to the diagonal residual representation,
whose characters are trivial on \(H\subset\Gamma^+\).
For every \(Z\subset B[\Gamma^+]\) and
\(k\in P_\ell^\Gamma(H)\), we claim
\[
 [k]-1\in
 \overline{(\mathrm{CH}(D)\cup Z)}^{\,\mathfrak m_B,\Gamma^+}
                 =D.\mathrm{chBar}(Z).              \tag{PQL.14}
\]
Fix a term of the defining intersection, with coefficient
power \(a\) and open normal subgroup \(V\subset\Gamma^+\),
and call its quotient algebra \(E\). The canonical
homomorphism \(\Gamma^+\to E^\times\) has open normal
kernel \(U\), because it kills \(V\).
For \(h\in H\), let \(x\) be its image and \(p\) its
\(D\)-characteristic polynomial. Cayley--Hamilton gives
\(p(x)=0\), and the residual identity gives
\[
 p(X)\equiv(X-1)^n\pmod{\mathfrak m_B}.              \tag{PQL.15}
\]
The coefficient ideal \(\mathfrak m_BE\) is nilpotent,
and \(\ell\in\mathfrak m_B\). In the commutative subalgebra
generated by \(x\) and the central coefficient ring,
(PQL.15) makes \(x-1\) nilpotent modulo a nilpotent
ideal, hence nilpotent. The commuting elements
\(\ell,x-1\) are nilpotent, so their generated ideal \(J_x\)
is nilpotent. For \(y\in J_x^b\), \(b\geq1\),
\[
 (1+y)^\ell-1\in J_x^{b+1},                         \tag{PQL.16}
\]
since \(\ell y\in J_x^{b+1}\) and \(y^j\in J_x^{b+1}\)
for every \(j\geq2\). Iteration starting at \(y=x-1\)
shows \(x^{\ell^e}=1\) for some \(e\).
Thus the image of \(H\) in the finite group
\(\Gamma^+/U\) is an \(\ell\)-group. PQL3, with
\(K=\Gamma^+\), puts \(k\) in \(U\).
This proves (PQL.14) in every quotient, including the
zero quotient, hence in their intersection.

For \(Z=\varnothing\), PQL8 gives
\[
 k\in\ker g\quad\Longrightarrow\quad [k]-1\in\ker D. \tag{PQL.17}
\]
For \(Z=\mathrm{pairGens}(u)\), (PQL.14) also places these
elements in the exact closed ideal used by
\(\mathrm{augmentationIdealKernel}\).
That second conclusion will be needed for nonsemisimplicity.
% END PQL9

% BEGIN PQL10
\paragraph{PQL10. Descent of the full determinant pair.}\label{claim:PQL10}
The kernel of
\(p_*:B[\Gamma^+]\to B[(\Gamma')^+]\) is generated by
\([k]-1\), \(k\in\ker g\). Choose representatives \(s(b)\)
of the quotient group. The difference
\([\gamma]-[s(g\gamma)]\) is a left multiple of
\([s(g\gamma)^{-1}\gamma]-1\); subtracting representatives
in a finite sum proves the assertion. By (PQL.17) and
PQL7 there is a unique full determinant \(D'\) with
\[
 D'\circ p_*=D,\qquad \psi'=\psi.                   \tag{PQL.18}
\]
This is equality of laws after arbitrary scalar extension.
All characteristic polynomials, coefficients and traces
therefore commute with pullback and coefficient change.

Use PQL6's residual datum \(r'\).
Continuity of the descended coefficient functions follows
from the quotient topology. The residual determinant
identity holds by surjective pullback. The local
character data remain unchanged. Polarization commutes
with \(p_*\), as does each ordered ordinary element
\[
 \prod_{i=1}^n
       \bigl([\phi_v(a_i)]-\psi_{v,i}(a_i)\bigr).    \tag{PQL.19}
\]
Lifting the test group element proves the downstairs
polarization and ordinary trace clauses.
For an inertia element downstairs, lift its original
local parameter \(a\in L.J_u\). Membership in \(r'.Ip_u\)
says that every residual character is one on its image.
Residual compatibility says the same upstairs, so the
original inertia clause applies. This checks every field
of \(\mathrm{Pair}\,r'\,B\), with the same ring \(B\),
\(O\)-algebra structure, ideal \(J\), and local characters.
% END PQL10

% BEGIN PQL11
\paragraph{PQL11. Every original point is unipotent on \(H\) at finite order.}
\label{claim:PQL11}
Let \(P:\mathrm{Pt}(T,u,d)\) be an original point,
\(A=K[[t]]\), and \(f=P.f:B\to A\).
The ring map \(B\to K\) obtained by taking constant
coefficients kills \(\mathfrak m_B\). Indeed its image is
a finite integral domain, hence a field; its kernel is
maximal, and \(B\) has just one maximal ideal.
No surjectivity of \(O\to K\) is needed here.
The finite field \(K\), being a \(\mathbb Z_\ell\)-algebra,
has characteristic \(\ell\). One may see this by noting
that the map from the infinite DVR \(\mathbb Z_\ell\)
to \(K\) has nonzero prime kernel, necessarily \((\ell)\).

For \(h\in H\subset\Gamma_0^+\), the full determinant
identity of \(P.\mathrm{maps}\) and the residual characters
give
\[
 \det(X-\rho(h))\bmod t=(X-1)^n.                    \tag{PQL.20}
\]
Cayley--Hamilton implies \((\rho(h)-1)^n\in
t\operatorname{Mat}_n(A)\). Thus
\((\rho(h)-1)^{nm}\in t^m\operatorname{Mat}_n(A)\).
In characteristic \(\ell\), commuting binomial expansion
gives
\(\rho(h)^{\ell^e}-1=(\rho(h)-1)^{\ell^e}\).
Taking any \(e\) with \(\ell^e\geq nm\) proves
\[
 \forall m\ \forall h\in H\ \exists e:\quad
      \rho(h^{\ell^e})\equiv1\pmod{t^m}.           \tag{PQL.21}
\]
This is exactly \(\mathrm{NP.UnipOn}(H)\). The proof
works for the entire matrix, not only the distinguished
block. It retains the original point and finite field.
% END PQL11

% BEGIN PQL12
\paragraph{PQL12. Simultaneous descent of a representation and an adjoint cocycle.}
\label{claim:PQL12}
Let \(R\) be a commutative ring of characteristic \(\ell\),
\(I\subset R\), and let \(\rho:\Gamma^+\to
\operatorname{GL}_n(R)\) and \(z:\Gamma^+\to\operatorname{Mat}_n(R)\)
be \(I\)-adically continuous with
\[
 z(ab)=z(a)+\rho(a)z(b)\rho(a)^{-1}.                 \tag{PQL.22}
\]
Assume the analogue of (PQL.21) modulo every \(I^m\).
Modulo \(I^m\), the subgroup
\[
 U_m=\{\gamma:\rho(\gamma)=1,\ z(\gamma)=0\}
\]
is open normal. Openness follows from local constancy
of the finitely many matrix entries. It is the kernel
of the homomorphism to the semidirect product
\(\operatorname{Mat}_n(R/I^m)\rtimes
\operatorname{GL}_n(R/I^m)\), so it is normal.
If \(\rho(a)=1\) modulo \(I^m\), (PQL.22) gives
\(z(a^\ell)=\ell z(a)=0\) modulo \(I^m\).
An extra factor \(\ell\) in the exponent of (PQL.21)
therefore puts each \(h\in H\) into \(U_m\).
PQL3 shows that every \(k\in P_\ell^\Gamma(H)\) satisfies
\[
 \rho(k)\equiv1\pmod{I^m},\qquad z(k)\equiv0\pmod{I^m}.
                                                            \tag{PQL.23}
\]
The argument uses no finiteness assumption on \(R/I^m\):
the subgroup \(U_m\) is open in a compact group, which
already makes its group quotient finite.

If \(R\) is \(I\)-adically separated, (PQL.23) gives
\(\rho(k)=1,z(k)=0\) exactly.
The representation factors through \(g\). The cocycle
does too: if \(g(a)=g(b)\), write \(b=a(a^{-1}b)\);
(PQL.22) and vanishing on the kernel give \(z(b)=z(a)\).
Continuity of the descended entries follows from the
quotient map. This is the actual content of
\texttt{coc\_pl} and the representation part of
\texttt{np\_desc}.
% END PQL12

% BEGIN PQL13
\paragraph{PQL13. The full \(\mathrm{NP}\) and \(\mathrm{Kap}\) records.}
\label{claim:PQL13}
Apply PQL12 with \(R=A=K[[t]]\), \(I=(t)\), first
with \(z=0\), then with each original Kap cocycle
\(\kappa_j\). Separation of \(A\) follows directly
from its coefficients. We obtain unique functions with
\[
 \rho'(g_0a)=\rho(a),\qquad
 \kappa'_j(g_0a)=\kappa_j(a).                       \tag{PQL.24}
\]
Construct \(x'=\mathrm{NP.desc}\) with the same \(K,A,t\),
block size \(a\), local matrices \(w,w^{-1}\), types,
ordinary characters \(\psi\), and all integral local data.
The polarization and block equations descend by lifting
a group element. The spanning condition remains true:
every original spanning matrix occurs as a descended
matrix at \(g_0(a)\), and preservation of \(\varepsilon\)
preserves the required subgroup \(G'\).
Inertia, triangularity, and the exact \(N1\) condition
use the same local parameters through \(g_0\phi_v=\phi'_v\).
This verifies every \(\mathrm{NP}\) field.

Set \(k'.g=k.g\) and use the descended \(\kappa'_j\);
keep each witness matrix \(Y_{j,v}\) unchanged.
For a continuous surjection preserving the global frame,
a matrix-valued function downstairs satisfies any one of
the source's Selmer conditions if and only if its pullback
does. This can be checked clause by clause: matrix-subspace
membership and the cocycle identity use surjectivity;
adic continuity uses the quotient topology; polarization
uses \(g_0(cac^{-1})=c'g_0(a)(c')^{-1}\); inertia and the
local coboundary witnesses use the unchanged local maps.
The same holds for a coboundary \(dY\), with exactly the
same matrix \(Y\) and its prescribed subspace and symmetry.

Consequently the cocycles and local witnesses satisfy
the Kap membership and local fields. A downstairs
linear dependence modulo boundaries pulls back to
one upstairs, giving the original independence.
Conversely pull back any downstairs Selmer cocycle,
apply the original Kap generation clause, and descend
the resulting coboundary identity by surjectivity.
The exponent of \(t\), coefficients, and witness matrix
are unchanged. This proves the complete \(\mathrm{Kap.desc}\),
including its generation field, without replacing Kap
by an abstractly isomorphic module.
% END PQL13

% BEGIN PQL14
\paragraph{PQL14. Descent of truncated Selmer cocycles.}\label{claim:PQL14}
At level \(N\), the quotient kernel is
\(P_\ell^{\Gamma_N}(h_N^{-1}H)\).
For \(A_m=A/(t^m)\), any
\(z\in\mathrm{SelP}(h_N,x,A_m)\) is a continuous
adjoint cocycle for \(\rho\circ h_N\).
The representation is of \(\ell\)-power order on \(h_N^{-1}H\)
by (PQL.21). PQL12, applied with a nilpotent ideal
(or directly in \(A_m\)), makes \(z\) zero on the
quotient kernel. It thus descends through \(g_N\).

For \(\mathrm{Sel1}\) and \(\mathrm{Sel2}\), the coefficient
ring is the source's
\[
 \mathcal D(A_m,g)=
 A_m\oplus\bigoplus_{j=1}^g A_m e_j
       \oplus\bigoplus_{j=1}^g A_m e'_j\oplus A_m\eta ,
\]
with its original multiplication. The only identity needed
here is that the ideal spanned by the \(e_j\) has square zero.
The representation defining these cocycles is
\[
 \rho_1(a)=\bigl(1+E(a)\bigr)\rho(a),\qquad
 E(a)=\sum_j e_j\kappa_j(a).                         \tag{PQL.25}
\]
It is a representation because the \(\kappa_j\) satisfy
(PQL.22) and \(e_ie_j=0\).
If \(b=h^{\ell^e}\) makes \(\rho(b)=1\) in \(A_m\),
then \(\rho_1(b)=1+E(b)\) and
\[
 \rho_1(b)^\ell=1+\ell E(b)=1.
\]
The coefficient ring has characteristic \(\ell\).
Thus \(\rho_1\) also has \(\ell\)-power order on the
chosen open subgroup. Apply PQL12 to each cocycle
over \(\mathcal D(A_m,g)\); the ideal generated by \(t\)
is nilpotent, so (PQL.23) is exact after finitely many
powers. Every Sel1 or Sel2 cocycle descends. Its local,
polarization, and matrix-subspace clauses descend by
PQL13's clausewise argument. This proves
\texttt{sel\_desc} for the actual truncated coefficients.
% END PQL14

% BEGIN PQL15
\paragraph{PQL15. The same constants in B1 and B2.}\label{claim:PQL15}
Fix \(N,m\). Pullback along \(g_N\) injects the downstairs
Sel1 set into the upstairs Sel1 set, and identifies their
boundary sets bijectively. For the latter assertion,
every upstairs boundary is \(dY\) for a matrix \(Y\)
in the unchanged prescribed subspace; the same \(Y\)
gives a downstairs boundary with that pullback.
Equations (PQL.8) and (PQL.24) identify the two
representations and first-order representations exactly.
Hence whenever the original Sel1 set is finite,
\[
 |\mathrm{Sel1}'_{N,m}|\leq|\mathrm{Sel1}_{N,m}|,\qquad
 |\mathrm{Bd1}'_{N,m}|=|\mathrm{Bd1}_{N,m}|.
                                                            \tag{PQL.26}
\]
Given the original B1 constant \(C\), use the same eventual
set of \(N\) for each \(m\). Since \(K\) and \(d\) are
unchanged, (PQL.26) proves
\[
 |\mathrm{Sel1}'_{N,m}|
 \leq|\mathrm{Bd1}'_{N,m}|\,|K|^{md+C}.
                                                            \tag{PQL.27}
\]
Finiteness is included. The choice of \(C\) remains
before every \(m\).

For B2, pull back the entire supplied array of downstairs
SelP cocycles. At an original good level \(N\), B2 supplies
coefficients \(c_i\in A\), at least one a unit, and a
Sel2 cocycle \(\Psi\), with
\[
 (\Psi(\gamma)_{ab}).e'_j
       =\left(t^C\sum_i c_i z_{ij}(\gamma)\right)_{ab}.
                                                            \tag{PQL.28}
\]
PQL14 descends \(\Psi\); PQL13 descends its Sel2
membership. Surjectivity of \(g_N\) descends (PQL.28)
with the identical coefficients and the identical \(C\).
No larger eventual threshold is required.
This proves both implications in \texttt{b\_desc}
for arbitrary original Kap, including \(g>0\).
It does not assert that the first-order Selmer space vanishes.
% END PQL15

% BEGIN PQL16
\paragraph{PQL16. Each point retains its map and prime kernel.}
\label{claim:PQL16}
For an original \(P:\mathrm{Pt}(T,u,d)\), take
\(x',k'\) from PQL13 and the bounds from PQL15.
Keep its eigenvalue array \(\alpha_{j,i}\).
Its characteristic-polynomial identity is unchanged
because \(\rho'(g_0\sigma_j)=\rho(\sigma_j)\);
the injectivity of each eigenvalue list is identical.
Keep also its \(O\)-algebra structure on \(K\) and,
literally, the original homomorphism
\[
 f'=f:B\longrightarrow K[[t]].                     \tag{PQL.29}
\]
To verify the determinant field of \(\mathrm{maps}\),
pull both candidate laws back along \(g_0\).
PQL10's full-law equality, (PQL.24), and the original
point identity make them equal; surjectivity of the
group algebra makes the laws downstairs equal.
The local character part is unchanged.
Thus
\[
 P':\mathrm{Pt}(Q.\mathrm{grp},u',d),\qquad
 \ker f'=\ker f,\qquad K'=K.                        \tag{PQL.30}
\]
Applying this independently to a finite list preserves
every point in that list and every equality of residue-field
cardinalities. There is no new coefficient extension.
% END PQL16

% BEGIN PQL17
\paragraph{PQL17. All automorphic patching data descend on the same modules.}
\label{claim:PQL17}
Let \(W:\mathrm{TWAut}(u,J,\sigma,D_0,T)\) be the
given original automorphic datum. Its base datum can
be reused with the same weight variables, map
\(\lambda\), module \(M\), actions, freeness, finiteness,
and support identity
\[
 \operatorname{Ann}_B(M)\subset\mathfrak q
              \quad\Longleftrightarrow\quad J\subset\mathfrak q
       \qquad(\mathfrak q\ {\rm prime}).             \tag{PQL.31}
\]
The weight identities remain valid because the local
characters \(\psi\) and their residual values are unchanged.

At each \(N\), the residual characters of \(W.A_N\)
are the pullbacks of the base characters. They are
therefore trivial on \(H_N\), so PQL6 and PQL10
descend its residual datum and full pair on the
same ring \(R_N\). Keep the maps \(s_N\), elements
\(U_{N,j,i}\), module \(M_N\), exponent array, rank,
and additive trivialization \(e_N\) verbatim.
All their defining equations still hold.

The target splitting algebra for the comparison maps
is unchanged: its defining characteristic polynomials
are equal by (PQL.8) and (PQL.18), with the same
residue lifts \(ut_{j,i}\).
After this literal identification reuse
\(\theta_N,\varphi_N,\psi_N\). The full determinant
compatibility of \(\theta_N\) is checked by pulling
back to \(\Gamma_N^+\), where it is the original one.
The coefficient-map identity, the images of \(U\),
semilinearity, and both composition identities involve
only the unchanged rings, maps, modules and scalar
\[
 f_N=\prod_j\prod_i\prod_{i'\ne i}
 (U_{j,i}-U_{j,i'})(U_{j,i}-Nq(N,j)U_{j,i'}).
                                                            \tag{PQL.32}
\]
They are therefore retained exactly.

Finally the ordinary annihilator relations involve
traces of a group element multiplied by \(q1El\)
or \(q2El\), and the original common exponent \(W.ex\).
Those ordered products commute with the group-algebra
quotient; lift the test group element and apply the
original relation on \(M_N\). This proves both
\(\mathrm{TWMaps.rel1}\) and \(\mathrm{rel2}\).
We have obtained a full
\[
 W':\mathrm{TWAut}(u',J,\sigma,D_0,Q.\mathrm{grp})
                                                            \tag{PQL.33}
\]
without modifying \(J\), any ring, any module, any
support equality, or any comparison map.
This is transport of supplied finite-level data; it
does not itself construct the later limiting patched
ring or its support theorem.
% END PQL17

% BEGIN PQL18
\paragraph{PQL18. The precise closed augmentation ideals under the quotient.}
\label{claim:PQL18}
Write \(I_u=D.\mathrm{chBar}(\mathrm{pairGens}(u))\)
as a two-sided ideal, and similarly \(I_{u'}\).
We need
\[
 [k]-1\in I_u\quad(k\in\ker g),\qquad
             p_*^{-1}(I_{u'})\subset I_u.            \tag{PQL.34}
\]
The first statement is PQL9. To prove the second, first
check that every generator of
\(\mathrm{CH}(D')\cup\mathrm{pairGens}(u')\)
lifts to the corresponding generating set upstairs.
For Cayley--Hamilton coefficients, lift the finite
tuple of group-algebra elements defining the generic
polynomial. Polynomial-law equality (PQL.18) makes
each coefficient map to the desired coefficient.
Ordinary generators lift through the unchanged local
maps and characters. For inertia generators, lift
the local parameter, and use residual compatibility
as in PQL10 to preserve membership in \(Ip\).
Polarized versions lift because \(p_*\) commutes
with the polarization operation.

Fix \(a\) and \(V\triangleleft_o\Gamma^+\), and set
\(V'=g(V)\), an open normal subgroup of \((\Gamma')^+\).
Each generator of the downstairs ideal with
\(\mathfrak m_B^a,V'\) lifts to the upstairs ideal
with \(\mathfrak m_B^a,V\): for the subgroup generators,
choose their preimages inside \(V\).
Since \(p_*\) is surjective, multiplication on either
side also lifts. If \(p_*(z)\) belongs to that downstairs
ideal, subtract a lift belonging to the upstairs
ideal. The difference is in \(\ker p_*\), generated
by the \([k]-1\); PQL9 puts all these in every
upstairs ideal. Hence \(z\) belongs to this upstairs
ideal. Intersecting over \(a,V\) proves (PQL.34).
The use of \(g(V)\), rather than an unrelated subgroup
downstairs, keeps the exact two-parameter closure.
% END PQL18

% BEGIN PQL19
\paragraph{PQL19. Nonsemisimplicity and discriminants at the original primes.}
\label{claim:PQL19}
Fix a field-valued map \(f:B\to\kappa\).
Let \(E\) be the quotient of \(\kappa[\Gamma^+]\)
by the two-sided ideal generated by \(f(I_u)\);
define \(E'\) analogously from \(I_{u'}\).
The first statement of (PQL.34) gives a surjective
homomorphism \(\kappa[(\Gamma')^+]\to E\).
The second says that it kills \(f(I_{u'})\):
lift each original coefficient-ring element of
\(I_{u'}\) along \(p_*\), and apply (PQL.34).
Thus there is a surjection
\[
                         E'\longrightarrow E.      \tag{PQL.35}
\]
A quotient of a semisimple ring is semisimple.
For example, its regular module is a quotient of
the original semisimple regular module; as the
kernel acts by zero, semisimplicity is the same
over the quotient ring. Contraposition of this
fact in (PQL.35), using the very same \(f\), gives
\[
 \mathrm{NonSS}(I_u,\mathrm{id},\mathfrak q)
 \ \Longrightarrow\
 \mathrm{NonSS}(I_{u'},\mathrm{id},\mathfrak q)
                  \quad\hbox{for every }\mathfrak q.         \tag{PQL.36}
\]
Only this direction is needed or claimed.

The discriminant predicate is equivalent in both
directions. Its witness is a pair of \(n^2\)-tuples
of group elements such that the matrix
\((\operatorname{tr}_D(g_i h_j))_{i,j}\)
has determinant outside \(\mathfrak q\).
Push the tuples down, or lift them up; (PQL.18)
gives the identical scalar matrix. Hence
\[
 \mathrm{DiscSt}(D,\mathfrak q)
       \quad\Longleftrightarrow\quad
 \mathrm{DiscSt}(D',\mathfrak q).                   \tag{PQL.37}
\]
Together (PQL.36)--(PQL.37) transport every required
\(\mathrm{StOrTh}\) predicate in the forward direction.
All primes here remain ideals of the same \(B\).
% END PQL19

% BEGIN PQL20
\paragraph{PQL20. The original implication \(\mathrm{CoreBC}\Rightarrow
\mathrm{Raw}\).}\label{claim:PQL20}
Assume the full original \(\mathrm{CoreBC}\) statement,
and take arbitrary data satisfying the original
\(\mathrm{Raw}\) premises. Construct \(Q\) by PQL5,
\(r',u'\) by PQL6 and PQL10, \(W'\) by PQL17,
and each point \(P'_i\) by PQL16.
The source's adaptation condition remains true by
PQL6, and \((Q.\mathrm{grp}).\mathrm{ProL}\) is
(PQL.10). Every characteristic coefficient of
the original determinant is a coefficient of
the descended determinant at the image element;
the local characters are the same. Thus
\(\mathrm{pairVals}(u)\subset\mathrm{pairVals}(u')\).
The original density condition
\(\mathrm{GenBy}_O(\mathrm{pairVals}(u))\)
implies the corresponding one for \(u'\).
The equal-cardinality premise for the fields of
the points is preserved literally.

Apply \(\mathrm{CoreBC}\) to these descended data.
Its prime list is
\((\ker P'_i.f)_i=(\ker P_i.f)_i\), by (PQL.29).
Given the original intervening primes \(e_i\),
their nonsemisimplicity passes to the descended
datum by (PQL.36); all inclusions in adjacent
point primes are unchanged. A starting prime
\(\mathfrak q_1\) with discriminant nonvanishing
retains that property by (PQL.37), together with
the same conditions \(J\subset\mathfrak q_1\)
and \(\mathfrak q_1\subset\ker P_0.f\).
For any final prime \(\mathfrak q\) below
\(\ker P_h.f\), its original \(\mathrm{StOrTh}\)
condition passes forward by PQL19.
The conclusion of \(\mathrm{CoreBC}\) is exactly
\(J\subset\mathfrak q\), in the original ring.
This proves the complete \(\mathrm{ChainConcl}\),
for all the original \(h\), and hence \(\mathrm{Raw}\).
The hypothesis \(\ell\ne2\) is retained exactly at
the \(\mathrm{CoreBC}\) invocation.
% END PQL20

% BEGIN PQL21
\paragraph{PQL21. Effect on the selected Count-II dependency cut.}
\label{claim:PQL21}
On the selected \(h=0,a=n=4,g=0,\ell=3\) data,
PQL20 closes the entire pro-\(\ell\) quotient and
descent step left in CRW21. It does so on the
original \(B,J,P.f,K\), full determinant pair,
coupled B1/B2 witness, and supplied \(\mathrm{TWAut}\).
Its scope is stronger than the selected one-point
case, but it has not produced \(\mathrm{CoreBC}\).

The S4 construction in the companion LIM chapter
can be applied to the actual \(Q.\mathrm{grp}\)
from PQL5, because its exact ProL premise has now
been proved. No alternative quotient group or
point is chosen at that interface.
The remaining arithmetic producers in the Raw
route are the coefficient normalization and added
Unit parameter leading from Core to CoreBC, and
the compatible limiting patched ring, maps and
modules together with S5's support, dimension and
fiber contracts. The original finite-level modules
in PQL17 are already supplied by TWAut; their
limiting patching and support are distinct tasks.

For Count II as a whole, ordinary/control and
separating arithmetic fibers, the original type
and neat-level producer, the four further Hida
endpoints, and all-level same-\(J\) HasAut remain.
This chapter changes none of the two original
rank-four Links, the fixed complex embedding,
the common outer coefficient threshold, full
CohP, LocCFT, N4 or local hypotheses.
It therefore closes a substantial internal
producer while leaving the selected graph's
compound \(\mathrm{COUNT\_II\_SUPPORT}\) leaf open.
% END PQL21
